\section{Канал}
\label{sec:channel}

Дельта Qwen3.8-27B весит 34\,844 байта на каждую новую позицию и 156.9~МБ рекуррентного состояния сверху. Порция в
тысячу токенов даёт около 192~МБ, из которых 157~МБ приходятся на состояние; по каналу в 1~Гбит/с она идёт около
полутора секунд, по каналу в 100~Мбит/с около пятнадцати. Состояние можно передавать в 16-битных числах, тогда оно
вдвое меньше, но такой режим нужно проверить на эквивалентность генерации, и в расчётах ниже мы показываем оба
варианта.

\textbf{Один ход.} Таблица~\ref{tab:turn} показывает, сколько ждёт агент на одной порции у клиента и через сервер.
Порция в 10 тысяч новых токенов на глубине 200 тысяч занимает у M3 Max 6.6 минуты; через канал в 100~Мбит/с гибрид
отдаёт её за 42 секунды, через гигабит за 5 секунд. Малая порция в 2 тысячи токенов выигрывает меньше, потому что
основную часть дельты составляет рекуррентное состояние, а его размер от порции не зависит.

\begin{table}[t]
\centering
\caption{Время одной порции на M3 Max и через гибрид: подъём копии сессии в память GPU, prefill на одной B300,
передача дельты, вставка и последний токен на клиенте. В скобках ускорение. Qwen3.8-27B Q4, кэш в 8 битах, состояние в 32 битах; время на клиенте по замеренной
кривой prefill.}
\label{tab:turn}
\footnotesize
\begin{tabular}{lrrrrr}
\toprule
Порция, глубина & Клиент & 100 Мбит/с & 300 Мбит/с & 1 Гбит/с & 10 Гбит/с \\
\midrule
2 тыс., 100 тыс. & 43 с & 18 с (2.3×) & 6.4 с (6.7×) & 2.2 с (20×) & 0.5 с (81×) \\
10 тыс., 100 тыс. & 3.6 мин & 41 с (5.3×) & 14 с (15×) & 5.0 с (44×) & 1.3 с (164×) \\
2 тыс., 200 тыс. & 78 с & 19 с (4.2×) & 6.5 с (12×) & 2.3 с (34×) & 0.7 с (115×) \\
10 тыс., 200 тыс. & 6.6 мин & 42 с (9.6×) & 15 с (27×) & 5.3 с (75×) & 1.7 с (233×) \\
\bottomrule
\end{tabular}
\end{table}

\textbf{Разовая загрузка.} Иногда контекст приходит целиком: агент загружает репозиторий или длинный журнал,
запускается новый субагент, сессия начинается заново после сжатия. Тогда клиент читает весь контекст с нуля
(таблица~\ref{tab:oneshot}). 200 тысяч токенов занимают у M3 Max 69 минут, а гибрид через гигабитный канал отдаёт их
за 71 секунду: 14 секунд считает сервер, остальное занимает передача. Здесь гибрид выигрывает во времени на любом канале из таблицы; трафика при этом около 7~ГБ на загрузку.

\begin{table}[t]
\centering
\caption{Время до первого токена при загрузке всего контекста сразу, M3 Max и гибрид. Та же модель и те же форматы,
что в таблице~\ref{tab:turn}.}
\label{tab:oneshot}
\footnotesize
\begin{tabular}{lrrrrr}
\toprule
Контекст & Клиент & 100 Мбит/с & 300 Мбит/с & 1 Гбит/с & 10 Гбит/с \\
\midrule
100 тыс. токенов & 21 мин & 5.0 мин & 1.7 мин & 35 с & 8.6 с \\
200 тыс. токенов & 69 мин & 9.7 мин & 3.4 мин & 71 с & 20 с \\
262 тыс. токенов & 1.9 ч & 12.7 мин & 4.5 мин & 95 с & 28 с \\
\bottomrule
\end{tabular}
\end{table}

\textbf{Пример вне программирования.} Аналитик загружает в агента стенограмму встречи на 70 тысяч токенов и просит
резюме на 3 тысячи токенов (таблица~\ref{tab:transcript}). M3 Max читает стенограмму 11.5 минуты и ещё 6.3 минуты
пишет резюме. Через гигабитный канал сервер отдаёт память модели за 24 секунды, и вся задача занимает 6.7 минуты
вместо 17.8. Оставшееся время уходит на генерацию, её гибрид не ускоряет. MacBook Air с 32~ГБ читает стенограмму сам
больше 20 минут, и гибрид сокращает задачу вдвое. Через API такое резюме стоит от 0.6 до 2 центов. Выигрыш гибрида
здесь во времени, а со своим сервером prefill стенограмма к тому же не покидает компанию.

\begin{table}[t]
\centering
\caption{Стенограмма на 70 тысяч токенов и резюме на 3 тысячи: время до первого токена и время всей задачи. M3 Max
по замеренным кривым; M5 Max и MacBook Air оценены по пиковой производительности и полосе памяти. Сервер читает
стенограмму 3.5 секунды. У Air под систему оставлено 8~ГБ; с запасом 12~ГБ, как у остальных клиентов, задача в его
память не помещается.}
\label{tab:transcript}
\footnotesize
\begin{tabular}{lrrrrrr}
\toprule
 & \multicolumn{3}{c}{До первого токена} & \multicolumn{3}{c}{Вся задача} \\
\cmidrule(lr){2-4}\cmidrule(lr){5-7}
Клиент & сам & 200~Мбит/с & 1~Гбит/с & сам & 200~Мбит/с & 1~Гбит/с \\
\midrule
M3 Max, 128~ГБ & 11.5 мин & 1.8 мин & 24 с & 17.8 мин & 8.1 мин & 6.7 мин \\
M5 Max, 64~ГБ, оценка & 2.7 мин & 1.8 мин & 24 с & 6.8 мин & 5.9 мин & 4.5 мин \\
MacBook Air M5, 32~ГБ, оценка & 21.6 мин & 1.8 мин & 25 с & 38.1 мин & 18.3 мин & 16.9 мин \\
\bottomrule
\end{tabular}
\end{table}

\textbf{Когда оффлоад выигрывает.} Клиент отдаёт порцию серверу, если чтение на сервере вместе с передачей дельты
быстрее, чем чтение на клиенте. Таблица~\ref{tab:offload} считает это правило на всём профиле вызовов для клиента M3
Max с замеренной кривой prefill и сервера с одной B300. Чем медленнее канал, тем крупнее должна быть порция, чтобы
передача окупилась, и тем меньше вызовов уходит на сервер.

\begin{table}[t]
\centering
\caption{Доля вызовов, где оффлоад быстрее локального чтения, и средняя экономия времени на вызов по всему профилю.
Клиент M3 Max (замеренная кривая prefill), сервер с одной B300 (оценка эффективной производительности).}
\label{tab:offload}
\footnotesize
\begin{tabular}{lrrrr}
\toprule
Канал & 100 Мбит/с & 300 Мбит/с & 1 Гбит/с & 2.5 Гбит/с \\
\midrule
Состояние в 32 битах: доля вызовов & 29.6\% & 45.6\% & 60.8\% & 68.7\% \\
\quad экономия на вызов, с & 19.3 & 24.2 & 26.5 & 27.1 \\
Состояние в 16 битах: доля вызовов & 38.7\% & 55.0\% & 67.0\% & 72.3\% \\
\quad экономия на вызов, с & 21.5 & 25.3 & 26.9 & 27.3 \\
\bottomrule
\end{tabular}
\end{table}

Средняя экономия растёт медленнее доли вызовов, потому что основное время сберегают крупные порции, а они выигрывают
уже на медленном канале. Даже при 100~Мбит/с оффлоад снимает 19 из 28 секунд среднего ожидания на вызов; гигабитный
канал добавляет ещё семь.

\textbf{Где стоит сервер.} Сервер может стоять в локальной сети офиса, в корпоративной сети компании, в колокации
или у арендованного поставщика. В локальной сети офиса клиент подключён к серверу проводом на гигабит и выше, и за
трафик никто не платит. Корпоративная сеть связывает офисы и удалённых сотрудников каналами между площадками и VPN, и
её скорость и цену задают эти каналы. В колокации и у поставщика клиент получает дельты через интернет, и тогда важны
три вещи: скорость подключения пользователя, цена исходящего трафика у поставщика и тариф самого пользователя,
безлимитный он или с оплатой за объём.

\textbf{Сеть пользователя.} Медианная скорость фиксированного подключения в мире 135~Мбит/с, в США 314~Мбит/с
(сентябрь 2026 года \cite{ookla2026}). Удалённый сотрудник на таком подключении получает от схемы меньше, чем сотрудник в
офисе с проводной сетью: при 100~Мбит/с на сервер уходит почти каждый третий вызов, при гигабите три из пяти. Wi-Fi и
VPN снижают полезную скорость, и для них оффлоад окупается только на крупных порциях.

\textbf{Трафик как статья расходов.} На нашем профиле клиент M3 Max с гигабитным каналом скачивает в среднем
130~МБ на вызов открытой модели, считая и вызовы, которые он читает сам. За восьмичасовой день он успевает обслужить
около 620 вызовов, и за 21 рабочий день это около 1.7~ТБ. По цене колокации около \$5 за терабайт это \$8 в месяц, а
по тарифам крупных облаков, от \$50 до \$90 за терабайт в зависимости от объёма \cite{awsdatatransfer}, от \$84 до
\$152, то есть в 6--10 раз больше, чем серверное вычисление и хранение сессий для того же пользователя (около \$15,
раздел~\ref{sec:economics}).
Пользователю с тарифом, где объём ограничен или оплачивается, эти терабайты тоже нужно учесть.
Поэтому схеме нужен дешёвый трафик. Его дают свой сервер в локальной сети или в колокации с
выгодным тарифом на канал либо поставщик, у которого трафик включён в цену: у Lambda исходящий трафик бесплатен
\cite{lambdapricing}, у выделенных GPU-серверов Hetzner он не ограничен и не тарифицируется \cite{hetznergpu}. Где
сервер приходится держать в крупном облаке с платным трафиком, выгоднее остаться либо в чисто облачном
инференсе, либо в чисто локальном.

\textbf{Вывод.} Гибриду нужен канал от 100~Мбит/с до клиента и дешёвый трафик у сервера. На таком канале большая
часть ожидания prefill уходит, а на гигабите крупная порция и загрузка всего контекста ускоряются в десятки раз. Где
трафик платный по тарифам крупных облаков, выгоднее остаться в облаке или работать локально.
